% ============================================================================
%  04-04: INTELIGENCIA DE NEGOCIOS — 10 CONSULTAS CON RESULTADOS REALES
% ============================================================================

\section{Consultas de inteligencia de negocios}

Las siguientes diez consultas responden a requerimientos analíticos del sistema TechParts ERP.
El código SQL se incluye tal como se encuentra en el archivo \texttt{anexos/consultas\_y\_pruebas\_bd2.sql} y los resultados tabulares corresponden a la salida real obtenida al ejecutarlas sobre la instancia de Supabase tras cargar los datos de demostración (13 pedidos).
Todos los importes se expresan en bolivianos (Bs.).

\subsection{Consulta 1: Productos más vendidos por unidades}

\textbf{Requerimiento:} Identificar los componentes con mayor volumen de ventas (pedidos en estado Completado o Entregado) para priorizar el reabastecimiento.

\lstinputlisting[language=SQL, firstline=154, lastline=168, caption={Q1: top de componentes por unidades e ingresos}, label=lst:q1]{anexos/consultas_y_pruebas_bd2.sql}

\begin{center}
\captionof{table}{Resultado de Q1: componentes más vendidos}
\label{tab:q1}
\begin{tablaAPA}
\begin{tabular}{L{4.6cm} L{3.2cm} L{1.9cm} R{1.5cm} R{2.7cm}}
\toprule
\textbf{Producto} & \textbf{Categoría} & \textbf{Marca} & \textbf{Unid.} & \textbf{Ingreso (Bs.)} \\
\midrule
Intel Core Ultra 9 285K & Procesadores & Intel & 5 & 32500.00 \\
Samsung 990 Pro 2TB NVMe & Almacenamiento & Samsung & 5 & 8000.00 \\
NVIDIA RTX 5080 Ti & Tarjetas de Video & NVIDIA & 4 & 78000.00 \\
G.SKILL Trident Z5 RGB 64GB DDR5 & Memoria RAM & G.Skill & 2 & 4800.00 \\
Logitech G Pro X Superlight 2 & Perifericos & Logitech & 2 & 2400.00 \\
ASUS ROG Maximus Z890 Hero & Tarjetas Madre & ASUS & 1 & 4500.00 \\
\bottomrule
\end{tabular}
\end{tablaAPA}
\fuenteAPA{Salida de la consulta Q1 en Supabase (PostgreSQL).}
\end{center}

\textbf{Interpretación:} El procesador Intel Core Ultra 9 285K y el SSD Samsung 990 Pro lideran en unidades (5 cada uno), pero la RTX 5080 Ti, con solo 4 unidades, aporta Bs.\ 78000.00, es decir el 59.9\,\% de los Bs.\ 130200.00 que suman las seis filas.
La estrategia de inventario debe diferenciar entre productos de rotación alta y bajo valor unitario (almacenamiento, memoria) y productos de alto valor unitario (GPU), cuyo capital inmovilizado es mayor.

\subsection{Consulta 2: Ingresos mensuales por categoría}

\textbf{Requerimiento:} Visualizar la evolución mensual de ingresos por categoría de componente.

\lstinputlisting[language=SQL, firstline=171, lastline=182, caption={Q2: ingresos mensuales por categoría}, label=lst:q2]{anexos/consultas_y_pruebas_bd2.sql}

\begin{center}
\captionof{table}{Resultado de Q2: ingresos mensuales por categoría}
\label{tab:q2}
\begin{tablaAPA}
\begin{tabular}{L{2.0cm} L{4.2cm} R{1.8cm} R{3.0cm}}
\toprule
\textbf{Mes} & \textbf{Categoría} & \textbf{Unidades} & \textbf{Ingreso (Bs.)} \\
\midrule
2026-11 & Tarjetas de Video & 2 & 39000.00 \\
2026-11 & Procesadores & 2 & 13000.00 \\
2026-11 & Almacenamiento & 2 & 3200.00 \\
2026-10 & Tarjetas de Video & 1 & 19500.00 \\
2026-10 & Memoria RAM & 1 & 2400.00 \\
2026-10 & Perifericos & 2 & 2400.00 \\
2026-09 & Procesadores & 3 & 19500.00 \\
2026-09 & Tarjetas de Video & 1 & 19500.00 \\
2026-09 & Almacenamiento & 3 & 4800.00 \\
2026-09 & Tarjetas Madre & 1 & 4500.00 \\
2026-09 & Memoria RAM & 1 & 2400.00 \\
\bottomrule
\end{tabular}
\end{tablaAPA}
\fuenteAPA{Salida de la consulta Q2 en Supabase (PostgreSQL).}
\end{center}

\textbf{Interpretación:} Las tarjetas de video son la categoría de mayor ingreso en los tres meses (Bs.\ 19500.00 en septiembre y octubre y Bs.\ 39000.00 en noviembre).
Los totales por mes de Q2 (septiembre Bs.\ 50700.00, octubre Bs.\ 24300.00, noviembre Bs.\ 55200.00) solo consideran líneas con \texttt{ComponenteId}; la vista \texttt{ventas.vw\_ventas\_mensuales} reporta subtotales mayores en octubre (Bs.\ 61100.00) y noviembre (Bs.\ 65700.00), y la diferencia (Bs.\ 36800.00 y Bs.\ 10500.00) coincide con el precio de los combos \textit{Gaming PC Ultimate} y \textit{Setup Gaming Completo}, vendidos en esos meses.

\subsection{Consulta 3: Rentabilidad de combos}

\textbf{Requerimiento:} Calcular el ahorro efectivo de cada combo respecto del precio de sus componentes individuales.

\lstinputlisting[language=SQL, firstline=185, lastline=185, caption={Q3: rentabilidad de combos (vista analítica)}, label=lst:q3]{anexos/consultas_y_pruebas_bd2.sql}

\begin{center}
\captionof{table}{Resultado de Q3: ahorro efectivo por combo}
\label{tab:q3}
\begin{tablaAPA}
\begin{tabular}{L{4.2cm} R{2.5cm} R{2.8cm} R{1.9cm} R{2.4cm}}
\toprule
\textbf{Combo} & \textbf{Precio venta} & \textbf{Suma componentes} & \textbf{Desc.\ (\%)} & \textbf{Ahorro (Bs.)} \\
\midrule
Gaming PC Ultimate & 36800.00 & 38800.00 & 5.15 & 2000.00 \\
Starter Build DDR5 & 14250.00 & 15000.00 & 5.00 & 750.00 \\
Setup Gaming Completo & 10500.00 & 11200.00 & 6.25 & 700.00 \\
\bottomrule
\end{tabular}
\end{tablaAPA}
\fuenteAPA{Salida de la consulta Q3 (vista \texttt{ventas.vw\_combos\_rentabilidad}) en Supabase.}
\end{center}

\textbf{Interpretación:} El descuento porcentual almacenado en \texttt{catalogo.Combos} coincide con el ahorro calculado a partir de los precios de los componentes (por ejemplo, 2000.00 sobre 38800.00 equivale a 5.15\,\%), lo que confirma la consistencia entre el precio de venta del combo y su composición en \texttt{ComboComponentes}.
\textit{Gaming PC Ultimate} ofrece el mayor ahorro absoluto (Bs.\ 2000.00) y \textit{Setup Gaming Completo} el mayor porcentaje (6.25\,\%).

\subsection{Consulta 4: Clientes con mayor gasto acumulado}

\textbf{Requerimiento:} Identificar a los clientes con mayor gasto total para diseñar programas de fidelización.

\lstinputlisting[language=SQL, firstline=188, lastline=188, caption={Q4: top de clientes mediante función almacenada}, label=lst:q4]{anexos/consultas_y_pruebas_bd2.sql}

\begin{center}
\captionof{table}{Resultado de Q4: clientes por gasto acumulado}
\label{tab:q4}
\begin{tablaAPA}
\begin{tabular}{L{3.6cm} L{5.4cm} R{1.7cm} R{2.6cm}}
\toprule
\textbf{Cliente} & \textbf{Correo} & \textbf{Pedidos} & \textbf{Gasto (Bs.)} \\
\midrule
Ana García Pérez & \texttt{cliente@techparts.com} & 9 & 127716.00 \\
Carlos Mendoza Rivera & \texttt{gerente.general@techparts.com} & 1 & 52780.00 \\
Laura Sánchez Morales & \texttt{vendedor@techparts.com} & 1 & 25404.00 \\
\bottomrule
\end{tabular}
\end{tablaAPA}
\fuenteAPA{Salida de \texttt{ventas.fn\_top\_clientes(10)} en Supabase.}
\end{center}

\textbf{Interpretación:} Un solo cliente concentra el 62.0\,\% del gasto acumulado (Bs.\ 127716.00 de Bs.\ 205900.00) y 9 de los 11 pedidos vigentes.
Esta concentración justifica un programa de fidelización escalonado y evidencia el riesgo de depender de pocas cuentas.

\subsection{Consulta 5: Abandono de carrito (JSONB)}

\textbf{Requerimiento:} Cuantificar el monto en riesgo por carritos no convertidos, consultando el contenido JSONB mediante el operador de contención \texttt{@>} (acelerado por el índice GIN).

\lstinputlisting[language=SQL, firstline=191, lastline=199, caption={Q5: carritos abandonados por mes con filtro JSONB}, label=lst:q5]{anexos/consultas_y_pruebas_bd2.sql}

\begin{center}
\captionof{table}{Resultado de Q5 y estado de la tabla \texttt{ventas.CarritosAbandonados}}
\label{tab:q5}
\begin{tablaAPA}
\begin{tabular}{L{4.4cm} R{1.2cm} R{2.6cm} R{1.6cm} L{2.6cm}}
\toprule
\textbf{Consulta / registro} & \textbf{Id} & \textbf{Total (Bs.)} & \textbf{Ítems} & \textbf{Convertido} \\
\midrule
Q5 (filtro \texttt{NOT Convertido}) & \multicolumn{4}{l}{0 filas devueltas} \\
Carrito registrado & 1 & 19500.00 & 1 & verdadero \\
\bottomrule
\end{tabular}
\end{tablaAPA}
\fuenteAPA{Salida de la consulta Q5 y consulta directa a \texttt{ventas.CarritosAbandonados} en Supabase.}
\end{center}

\textbf{Interpretación:} La consulta no devuelve filas porque el único carrito registrado (RTX 5080 Ti, Bs.\ 19500.00) figura como convertido: el procedimiento \texttt{sp\_confirmar\_venta} marca \texttt{Convertido = TRUE} para el usuario que concreta la compra, de modo que, tras la prueba transaccional con \texttt{COMMIT}, el carrito dejó de contarse como abandonado.
El resultado vacío es por tanto coherente con el diseño y muestra que la tabla se mantiene sincronizada con el flujo de ventas.

\subsection{Consulta 6: Alerta de stock crítico}

\textbf{Requerimiento:} Listar los componentes activos cuyo stock es igual o inferior al mínimo configurado, indicando su proveedor para emitir órdenes de compra.

\lstinputlisting[language=SQL, firstline=202, lastline=216, caption={Q6: alerta de stock crítico con proveedor}, label=lst:q6]{anexos/consultas_y_pruebas_bd2.sql}

\begin{center}
\captionof{table}{Resultado de Q6: componentes bajo el stock mínimo}
\label{tab:q6}
\begin{tablaAPA}
\footnotesize
\begin{tabular}{L{3.6cm} L{1.5cm} L{1.5cm} R{1.0cm} R{1.1cm} L{2.9cm} L{2.5cm}}
\toprule
\textbf{Componente} & \textbf{Categ.} & \textbf{Marca} & \textbf{Stock} & \textbf{Mín.} & \textbf{Proveedor} & \textbf{Contacto} \\
\midrule
Lian Li O11 Dynamic EVO & Cases & Lian Li & 4 & 5 & Cooler Master Latam & \texttt{cm@latam.com} \\
\bottomrule
\end{tabular}
\end{tablaAPA}
\fuenteAPA{Salida de la consulta Q6 en Supabase (PostgreSQL).}
\end{center}

\textbf{Interpretación:} Solo un componente de los 15 activos está bajo su mínimo: el gabinete Lian Li O11 Dynamic EVO (4 unidades frente a un mínimo de 5), cuyo proveedor registrado es Cooler Master Latam.
El campo \texttt{StockMinimo} por producto, incorporado en el rediseño, permite que la alerta sea específica en lugar de usar un umbral global.

\subsection{Consulta 7: Método de pago preferido}

\textbf{Requerimiento:} Determinar la distribución de métodos de pago en los pedidos no devueltos ni cancelados.

\lstinputlisting[language=SQL, firstline=219, lastline=232, caption={Q7: distribución de métodos de pago con función ventana}, label=lst:q7]{anexos/consultas_y_pruebas_bd2.sql}

\begin{center}
\captionof{table}{Resultado de Q7: distribución de métodos de pago}
\label{tab:q7}
\begin{tablaAPA}
\begin{tabular}{L{4.4cm} R{2.0cm} R{3.0cm} R{2.0cm}}
\toprule
\textbf{Método de pago} & \textbf{Pedidos} & \textbf{Monto (Bs.)} & \textbf{\%} \\
\midrule
Tarjeta & 6 & 120756.00 & 54.55 \\
QR / Billetera Digital & 3 & 29580.00 & 27.27 \\
Efectivo & 2 & 55564.00 & 18.18 \\
\bottomrule
\end{tabular}
\end{tablaAPA}
\fuenteAPA{Salida de la consulta Q7 en Supabase (PostgreSQL).}
\end{center}

\textbf{Interpretación:} La tarjeta concentra el 54.55\,\% de los pedidos, pero el efectivo registra el ticket promedio más alto (Bs.\ 27782.00 frente a Bs.\ 20126.00 de la tarjeta y Bs.\ 9860.00 del QR).
La transferencia no aparece porque todos sus pedidos se encuentran devueltos o cancelados y el filtro \texttt{Estado NOT IN (3, 4)} los excluye.
La suma de montos (Bs.\ 205900.00) coincide con el gasto total de la Tabla~\ref{tab:q4}, lo que valida la consistencia entre consultas.

\subsection{Consulta 8: Tiempo promedio de entrega por sucursal}

\textbf{Requerimiento:} Medir la eficiencia operativa de cada sucursal mediante el tiempo promedio entre pedido y entrega.

\lstinputlisting[language=SQL, firstline=235, lastline=246, caption={Q8: tiempo promedio de entrega por sucursal y región}, label=lst:q8]{anexos/consultas_y_pruebas_bd2.sql}

\begin{center}
\captionof{table}{Resultado de Q8: días promedio de entrega}
\label{tab:q8}
\begin{tablaAPA}
\begin{tabular}{L{2.6cm} L{3.6cm} L{2.4cm} R{2.4cm} R{2.0cm}}
\toprule
\textbf{Región} & \textbf{Sucursal} & \textbf{Ciudad} & \textbf{Entregados} & \textbf{Días prom.} \\
\midrule
Cochabamba & Sucursal Centro & Cochabamba & 6 & 1.83 \\
La Paz & Sucursal La Paz & La Paz & 1 & 3.00 \\
\bottomrule
\end{tabular}
\end{tablaAPA}
\fuenteAPA{Salida de la consulta Q8 en Supabase (PostgreSQL).}
\end{center}

\textbf{Interpretación:} La Sucursal Centro entrega en 1.83 días en promedio (6 pedidos) y la Sucursal La Paz en 3.00 días (1 pedido), es decir 1.17 días más.
Con un solo pedido entregado en La Paz la diferencia no es estadísticamente concluyente; la consulta queda disponible para reevaluarse cuando el volumen crezca.

\subsection{Consulta 9: Componentes sin ventas y mermas (ODS 12)}

\textbf{Requerimiento:} Detectar componentes sin ventas registradas, su estado de obsolescencia y las unidades dadas de baja en \texttt{auditoria.MermasInventario}.

\lstinputlisting[language=SQL, firstline=249, lastline=263, caption={Q9: componentes sin ventas y resumen de mermas}, label=lst:q9]{anexos/consultas_y_pruebas_bd2.sql}

\begin{center}
\captionof{table}{Resultado de Q9: componentes sin ventas}
\label{tab:q9}
\begin{tablaAPA}
\begin{tabular}{R{0.8cm} L{4.4cm} L{3.0cm} R{1.2cm} L{2.8cm} R{1.4cm}}
\toprule
\textbf{Id} & \textbf{Componente} & \textbf{Categoría} & \textbf{Stock} & \textbf{Obsolescencia} & \textbf{Merma} \\
\midrule
10 & Razer BlackWidow V4 Pro & Perifericos & 30 & Vigente & 0 \\
14 & Corsair MP700 PRO 2TB Gen5 & Almacenamiento & 25 & Vigente & 0 \\
9 & Cooler Master MasterLiquid 360 ATMOS & Refrigeracion & 22 & Vigente & 0 \\
6 & Corsair HX1500i Platinum & Fuentes de Poder & 20 & Vigente & 0 \\
11 & AMD Ryzen 9 9950X & Procesadores & 15 & Vigente & 0 \\
12 & ASUS ROG Crosshair X670E Hero & Tarjetas Madre & 12 & PorDescontinuar & 1 \\
13 & AMD Radeon RX 7900 XTX & Tarjetas de Video & 8 & Vigente & 0 \\
15 & Lian Li O11 Dynamic EVO & Cases & 4 & Vigente & 0 \\
\bottomrule
\end{tabular}
\end{tablaAPA}
\fuenteAPA{Salida de la consulta Q9 en Supabase (PostgreSQL).}
\end{center}

\textbf{Interpretación:} Ocho de los 15 componentes no registran ventas y mantienen 136 unidades en inventario.
La placa ASUS ROG Crosshair X670E Hero ya figura como \texttt{PorDescontinuar} y tiene una unidad dada de baja en el registro de mermas, lo que evidencia el seguimiento del ciclo de vida del hardware que respalda el ODS~12; los productos de este listado son candidatos a liquidación o revisión de catálogo para evitar obsolescencia.

\subsection{Consulta 10: Dashboard ejecutivo mensual}

\textbf{Requerimiento:} Generar un resumen mensual con indicadores clave para la gerencia general, validando además las funciones almacenadas.

\lstinputlisting[language=SQL, firstline=266, lastline=276, caption={Q10: indicadores ejecutivos mensuales}, label=lst:q10]{anexos/consultas_y_pruebas_bd2.sql}

\begin{center}
\captionof{table}{Resultado de Q10: indicadores ejecutivos mensuales}
\label{tab:q10}
\begin{tablaAPA}
\begin{tabular}{L{1.9cm} R{1.7cm} R{1.9cm} R{2.9cm} R{2.5cm} R{2.3cm}}
\toprule
\textbf{Mes} & \textbf{Pedidos} & \textbf{Clientes} & \textbf{Facturación (Bs.)} & \textbf{Ticket (Bs.)} & \textbf{Dev./Canc.} \\
\midrule
2026-11 & 4 & 2 & 76212.00 & 19053.00 & 0 \\
2026-10 & 5 & 2 & 93960.00 & 18792.00 & 2 \\
2026-09 & 4 & 1 & 58812.00 & 14703.00 & 0 \\
\bottomrule
\end{tabular}
\end{tablaAPA}
\fuenteAPA{Salida de la consulta Q10 en Supabase (PostgreSQL).}
\end{center}

La función \texttt{catalogo.fn\_recalcular\_stock\_combo} se validó con la consulta de las líneas 282--283 del script de pruebas, que compara el stock registrado de cada combo con el máximo que permite el inventario actual de sus componentes.

\begin{center}
\captionof{table}{Validación de \texttt{catalogo.fn\_recalcular\_stock\_combo}}
\label{tab:fn-stock-combo}
\begin{tablaAPA}
\begin{tabular}{L{4.6cm} R{3.0cm} R{4.0cm}}
\toprule
\textbf{Combo} & \textbf{Stock registrado} & \textbf{Máximo calculado} \\
\midrule
Gaming PC Ultimate & 5 & 10 \\
Setup Gaming Completo & 8 & 8 \\
Starter Build DDR5 & 8 & 10 \\
\bottomrule
\end{tabular}
\end{tablaAPA}
\fuenteAPA{Salida de la consulta de validación en Supabase (PostgreSQL).}
\end{center}

\textbf{Interpretación:} Septiembre, octubre y noviembre suman 13 pedidos y Bs.\ 228984.00; la diferencia de Bs.\ 23084.00 respecto del gasto acumulado de la Tabla~\ref{tab:q4} proviene de pedidos fuera de los estados Completado y Entregado (los dos pedidos devueltos o cancelados de octubre, 40\,\% de los pedidos de ese mes).
Respecto de los combos, el stock registrado de \textit{Gaming PC Ultimate} y \textit{Starter Build DDR5} es inferior al que permitiría el inventario de componentes (10 unidades), por lo que existe capacidad para reponer stock de combos sin comprar insumos adicionales.
